This chapter evaluates the proposed system. Because music generation quality
is inherently subjective, the evaluation combines \emph{automatic metrics}
(objective and reproducible) with \emph{qualitative analysis} (expert review
of sample outputs). The chapter reports quantitative results on a 20-sample
benchmark, analyzes fallback behavior and latency, compares L2M with related
systems, and discusses the findings, limitations, and ethical
considerations.

\section{Evaluation Methodology}
\label{sec:ev-methodology}

\subsection{Automatic Metrics}

Five automatic metrics are computed over all benchmark runs:

\begin{table}[ht]
\centering
\caption[Automatic evaluation metrics]{Automatic evaluation metrics and
their definitions.}
\label{tab:metric-definitions}
\begin{tabular}{p{4.4cm} p{7.7cm}}
\toprule
\textbf{Metric} & \textbf{Definition} \\
\midrule
Syllable--note alignment & Fraction of outputs where the number of generated
notes equals the total syllable count \\
Emotion--key consistency & Fraction of outputs using a key consistent with
the detected emotion per \ref{tab:emotion-mapping} \\
Tempo appropriateness & Fraction of outputs with tempo within the expected
BPM range for the detected emotion \\
MIDI validity & Fraction of MIDI files successfully re-parsed by music21
without errors \\
LLM success rate & Fraction of pipeline runs completed without activating a
fallback \\
\bottomrule
\end{tabular}
\end{table}

\subsection{Qualitative Assessment}

Three dimensions are evaluated by expert listening: \textbf{melodic
coherence} (smoothness of pitch progression, singability, phrase closure),
\textbf{emotional alignment} (perceived match between lyrical sentiment and
musical mood), and \textbf{rhythmic naturalness} (naturalness of note
duration patterns relative to speech rhythm).

\section{Test Dataset}
\label{sec:ev-dataset}

The benchmark comprises 20 lyrical inputs designed to span the emotion space
and a range of lyric complexity (\ref{tab:dataset}). Ten inputs use
literal, simple language and ten use poetic or metaphorical language.

\begin{table}[ht]
\centering
\caption[Benchmark composition]{Composition of the 20-sample evaluation
benchmark.}
\label{tab:dataset}
\begin{tabular}{l p{8.6cm}}
\toprule
\textbf{Dimension} & \textbf{Distribution} \\
\midrule
Emotion & Happy (5), Sad (5), Hopeful (4), Tense (3), Calm (2), Excited (1) \\
Length & Short (1 line, 5 samples), Medium (2--3 lines, 10 samples), Long
(4+ lines, 5 samples) \\
Style & Literal/simple (10), Poetic/metaphorical (10) \\
\bottomrule
\end{tabular}
\end{table}

Representative test cases include \enquote{The sun will rise again}
(hopeful, 6 syllables), \enquote{Dancing in the moonlight, feeling so
alive} (happy, 12 syllables), \enquote{Memories fade like photographs left
in the rain} (sad, 13 syllables), \enquote{Shadows creeping closer,
darkness all around} (tense, 11 syllables), and \enquote{Gentle waves upon
the shore, peaceful and serene} (calm, 12 syllables).

\section{Quantitative Results}
\label{sec:ev-quantitative}

\subsection{Automatic Metrics}

\ref{tab:results} reports the automatic metric scores across the
benchmark.

\begin{table}[ht]
\centering
\caption[Automatic metric results]{Automatic metric results on the
20-sample benchmark.}
\label{tab:results}
\begin{tabular}{l c p{5.6cm}}
\toprule
\textbf{Metric} & \textbf{Score} & \textbf{Details} \\
\midrule
Syllable--note alignment & \textbf{100\%} & All 20/20 outputs matched the
syllable count exactly \\
Emotion--key consistency & \textbf{95\%} & 19/20 used appropriate keys; one
edge case with ambiguous mixed emotion \\
Tempo appropriateness & \textbf{90\%} & 18/20 within expected BPM ranges;
two borderline outliers \\
MIDI validity & \textbf{100\%} & All 20/20 MIDI files parsed successfully by
music21 \\
LLM success rate & \textbf{85\%} & 17/20 runs used LLM output; 3/20
activated a fallback \\
\bottomrule
\end{tabular}
\end{table}

The perfect alignment score is a direct consequence of the hard
\enquote{EXACTLY $N$ notes} prompt constraint combined with schema
validation and, where necessary, fallback regeneration; alignment is the
single most important requirement for singability and was the primary
failure mode of unconstrained prompting (\Autoref{sec:ev-discussion}).

\subsection{Fallback Performance}

Of the three runs that activated a fallback, two were caused by network
timeouts on LLM API calls and one by a malformed JSON response (a note-count
mismatch). All three fallback-generated melodies passed the MIDI validity
check. Qualitative review rated the fallback outputs as \enquote{mechanically
correct but less expressive} than LLM outputs --- an expected trade-off of
deterministic generation, and an acceptable one given that the alternative
is outright failure.

\subsection{Latency Profile}

\ref{tab:latency} breaks down the mean end-to-end latency of
approximately 3.2 seconds. The two LLM calls dominate the runtime; notation
export and audio synthesis are comparatively negligible.

\begin{table}[ht]
\centering
\caption[Latency profile]{Mean per-operation latency across benchmark runs.}
\label{tab:latency}
\begin{tabular}{l c}
\toprule
\textbf{Operation} & \textbf{Mean time} \\
\midrule
Emotion analysis (LLM) & $\sim$1.4 s \\
Melody generation (LLM) & $\sim$1.1 s \\
MIDI/MusicXML export & $\sim$0.3 s \\
Audio rendering (WAV) & $\sim$0.4 s \\
\midrule
\textbf{End-to-end (with audio)} & \textbf{$\sim$3.2 s} \\
\bottomrule
\end{tabular}
\end{table}

\section{Qualitative Analysis of Sample Outputs}
\label{sec:ev-qualitative}

\subsection{Hopeful Sample}

For the input \enquote{The sun will rise again} (hopeful, 6 syllables), the
system detected the emotion \emph{hopeful} and selected G~major at 90~BPM in
4/4. The generated melody,
\begin{center}
G4\,(0.5) $\rightarrow$ A4\,(0.5) $\rightarrow$ B4\,(1.0) $\rightarrow$
C5\,(1.0) $\rightarrow$ B4\,(0.5) $\rightarrow$ A4\,(1.5),
\end{center}
forms an ascending arch that rises to the phrase peak (C5) and resolves
back down. All notes are diatonic to G~major and the entire melody spans a
singable fifth; the ascending contour is emotionally appropriate for the
hopeful register.

\subsection{Sad Sample}

For the input \enquote{Memories fade like photographs left in the rain}
(sad, 13 syllables), the system detected \emph{sad} and selected A~minor at
65~BPM in 4/4. The melody descends stepwise (A4 $\rightarrow$ G4
$\rightarrow$ F4 $\rightarrow$ E4 $\rightarrow$ D4 $\ldots$), resolving on
D4. The stepwise descent and slow tempo capture the melancholic mood
effectively, and the melody avoids any large major-interval leaps that would
be inconsistent with the emotional register.

\section{Comparison with Related Systems}
\label{sec:ev-comparison}

\ref{tab:comparison} positions L2M against representative supervised
lyrics-to-melody systems.

\begin{table}[ht]
\centering
\caption[Comparison with related systems]{Comparison of L2M with
representative lyrics-to-melody systems.}
\label{tab:comparison}
\begin{tabular}{l c c c c}
\toprule
\textbf{System} & \textbf{Training} & \textbf{Output formats} &
\textbf{Fallback} & \textbf{Open source} \\
\midrule
Bao et al.~\cite{bao19neural} & Yes (large corpus) & MIDI & No & No \\
ReLyMe~\cite{zhang22relyme} & Yes & MIDI & Partial & No \\
SongComposer~\cite{ding24songcomposer} & Yes (paired data) & MIDI + lyrics &
No & No \\
\textbf{L2M (ours)} & \textbf{No} & \textbf{MIDI, MusicXML, WAV, MP3} &
\textbf{Yes} & \textbf{Yes} \\
\bottomrule
\end{tabular}
\end{table}

L2M's advantages over supervised systems are the absence of any training
data or GPU requirement, a 100\% completion rate through the hybrid
architecture, multiple output formats including playable audio, and
sub-5-second end-to-end latency. Its limitations relative to fine-tuned
systems are less sophisticated melodic structures (single voice, no
harmony), limited style diversity, and musical creativity bounded by the
LLM's zero-shot capability.

\section{Discussion}
\label{sec:ev-discussion}

Four findings emerge from the evaluation, answering the research questions
posed in \Autoref{chap:intro}.

\textbf{Finding 1 --- LLM viability for cross-modal generation (RQ1).}
GPT-4o-mini demonstrates sufficient musical knowledge to generate
syllabically aligned, emotionally consistent melodies in a zero-shot
setting. This suggests that LLM pre-training encodes substantial implicit
music theory knowledge that structured prompting can elicit.

\textbf{Finding 2 --- Prompt engineering as a core technical contribution
(RQ4).} The syllable-count hard constraint proved critical: without explicit
enforcement, early prompt versions produced note counts deviating by
$\pm30\%$ from the syllable count. Structured JSON schemas and few-shot
examples reduced malformed response rates from roughly 35\% (unstructured
prompts) to roughly 15\%.

\textbf{Finding 3 --- Hybrid architecture value (RQ3).} The fallback
activation rate of 15\% (3/20 runs) demonstrates that LLM-only pipelines are
insufficient for production systems requiring reliability. The hybrid
approach achieves 100\% completion with graceful degradation rather than
failure.

\textbf{Finding 4 --- Two-stage decomposition (RQ2).} Separating emotion
analysis from melody generation yields two benefits: \emph{interpretability}
--- users can inspect and override the emotion reading before melody
generation --- and \emph{targeted fallback} --- if only melody generation
fails, the correct emotional context is preserved for the fallback
generator.

\section{Limitations}
\label{sec:ev-limitations}

Five limitations bound the present work:

\begin{enumerate}
    \item \textbf{Harmonic depth.} Generated melodies are monophonic; chord
    progressions, harmonization, and countermelody are outside the current
    scope.
    \item \textbf{Cultural scope.} The emotion--key mapping and contour
    algorithms reflect Western tonal music theory; non-Western scales,
    modes, and rhythmic systems are not supported.
    \item \textbf{Long-form coherence.} The chunking strategy maintains
    local melodic continuity via three-note context carryover but does not
    enforce global structural coherence (e.g., return to the tonic at
    phrase boundaries).
    \item \textbf{Evaluation scale.} The 20-sample benchmark enables
    proof-of-concept validation but is insufficient for statistically robust
    claims about generalization.
    \item \textbf{LLM API dependency.} The primary generation path requires
    OpenAI API access, incurring per-request cost and network latency.
\end{enumerate}

\section{Ethical Considerations}
\label{sec:ev-ethics}

Three ethical aspects deserve note. \textbf{Copyright}: generated melodies
may inadvertently resemble copyrighted works, and users should exercise
discretion before commercial use. \textbf{Attribution}: AI-generated content
should be clearly disclosed, particularly in creative or academic contexts.
\textbf{Access equity}: API costs may create barriers for users in
resource-constrained settings; support for locally hosted open-source LLMs
would mitigate this and is identified as future work.
